\section{Construction and rank of the boundary period lattice}
\label{sec:boundary-periods}

Let \(T\) be the complete discrete valuation ring on a horizontal
boundary trait, \(K=\operatorname{Frac}T\), \(\pi\) a uniformizer, and
\(G/T\) the semiabelian model with principally polarized abelian
generic fiber. At a nonsplit point the source first splits the toric
part of \(G_{\kappa(T)}\). In the horizontal case its finite splitting
field is separable because \(\kappa(T)\) has characteristic zero.
Lifting this extension to the finite unramified extension of \(T\)
gives the split special-fiber torus. Write \(T\) for the resulting
ring. When a split presentation has already been supplied, the
construction below starts there. The torus need not have a
complement in the special fiber.

\subsection{The formal torus and the two lattices}

Put \(S_e=\operatorname{Spec}(T/\pi^{e+1})\). The torus with character
group \(X\) has over \(S_0\) a kernel embedding into \(G_0\), with
abelian quotient \(B_0\). On every thickening there is an extension
\[
  1\longrightarrow \operatorname{Hom}(X,\mathbb G_m)_{S_e}
  \longrightarrow G_e\longrightarrow B_e\longrightarrow 0.
\]
Here one uses the formal Raynaud extension theorem for a smooth
semiabelian scheme whose special torus is isotrivial. In the present
case that torus is split. The morphisms commute with restriction to
each preceding thickening. The construction used in the release starts
with an even trivialization of the cubical line on the special torus,
lifts that trivialization and the torus kernel across each nilpotent
thickening, and descends the line to \(M_e\) on \(B_e\). Restricting a
descended cubical line back to the preceding level gives its previous
descent. The finite-level quotient is a torus torsor, so the
weight-\(x\) functions on it form an invertible sheaf. Multiplication
identifies the tensor product of the weight-\(x\) and weight-\(x'\)
sheaves with the weight-\(x+x'\) sheaf, as one sees on any
trivialization of the torsor. This supplies the homomorphism into the
dual abelian scheme on every level.

The rings \(T/\pi^{e+1}\) exhaust the formal completion, each \(B_e\)
is proper, and the compatible \(M_e\) have ample special fiber.
Grothendieck algebraization therefore gives a projective \(B/T\)
with an invertible sheaf \(M\) inducing these data. Full faithfulness
on proper schemes applied to the group operations and their identities
gives the group law on \(B\). Smoothness and geometric connectedness
yield the abelian scheme. The Raynaud extension \(G^\natural\) is an
extension of this \(B\) by the split torus. Its formal completion is
identified with the formal completion of \(G\). This accounts for the
presence of a torus at every infinitesimal level even though \(G_K\)
is abelian.
\footnote{\href{https://www.kwlan.org/articles/cpt-PEL-type-thesis-revision-dense.pdf}{K.-W.~Lan,
\emph{Arithmetic Compactifications of PEL-Type Shimura Varieties},
\S3.3.3, equations (3.3.3.1)--(3.3.3.2) and Proposition 3.3.3.6};
\href{https://stacks.math.columbia.edu/tag/089A}{Grothendieck
algebraization, Stacks Project, Theorem 30.28.4}.}

In the usual degeneration datum \(X\) is the character group, \(Y\)
the lattice of periods, and \(\phi:Y\to X\) the map attached to the
polarization. Lan's degree formula is
\[
  \deg\lambda_K=[X:\phi Y]^2\,\deg\lambda_B.
\]
A principal polarization on \(G_K\) therefore identifies \(Y=X\)
through \(\phi\), and \(\lambda_B\) identifies \(B\) with its dual.
This explains why one and the same lattice and a map \(c:X\to B(T)\)
occur in the local jet formulas. More intrinsically, \(c\) corresponds
under this principal identification to the preceding homomorphism
from weights into \(B^\vee(T)\). Its value is determined by the
rigidified weight bundle, so tensor product gives
\(c(x+x')=c(x)+c(x')\).
\footnote{\href{https://www.kwlan.org/articles/cpt-PEL-type-thesis-revision-dense.pdf}{Lan,
Lemma 3.4.4.3 and Theorem 4.2.1.14}. The hypotheses on the base in
\S4.1 hold with the complete normal noetherian domain \(T\) and its
radical maximal ideal. The model is semiabelian, with split special
torus and abelian generic fiber. The rigidified symmetric cubical line
is ample by Lemma 4.2.1.6. Its formal descent is the compatible system
of descended lines above, as required by Definition 4.2.1.1.}

\subsection{The valuation of the period factors}

Write \(f_x\) for the Fourier coefficients of a nonzero generator of
the level-one theta space. In the principal case the space is
one-dimensional. Every coefficient map is nonzero, and translating
the \(x\)-coefficient through a period \(y\) gives
\[
 f_{x+y}=\psi(y)\tau(y,x)\,t_{c(y)}^* f_x .
\]
The equality uses the canonical isomorphisms of the corresponding
line fibers, with the scalar convention of the displayed equation.
Lan proves nonvanishing and this formula in
\S\S4.3.2--4.3.3. In the release the nonvanishing step uses the
degree-two addition formula and torus-torsion weights. Its support is
a coset of a subgroup of odd index; the exclusion of prime-index
subgroup obstructions makes that subgroup the full lattice.

Let \(Q(y)=v_T(\psi(y))\), evaluated relative to an integral generator
of its one-dimensional line. Any two such generators differ by a
unit, so \(Q(y)\in\mathbb Z\) is unambiguous. Write \(B(y,x)\) for
\(v_T(\tau(y,x))\) relative to the integral biextension. Comparison of
the two orders of translating \(f_0\) gives
\[
  Q(y+z)=Q(y)+Q(z)+B(y,z),\qquad B(y,z)=B(z,y).
\]
The biextension law gives additivity of \(B\) in each argument.
Rigidification gives \(Q(0)=0\). The even cubical normalization of
the chosen line gives \(Q(-y)=Q(y)\). Taking \(z=-y\) yields
\(B(y,y)=2Q(y)\); by induction \(Q(ny)=n^2Q(y)\) for \(n\in\mathbb Z\).

Here is the convergence input giving strict positivity. Since \(f_0\)
is a nonzero integral section on the proper abelian quotient, some
germ has order at most \(m\) for a fixed \(m\). Indeed a section
vanishing at every finite level is zero by formal separation.
Translation by the integral point \(c(y)\) moves this witness to a
germ where \(t_{c(y)}^*f_0\) has order at most the same \(m\). At any
fixed finite level a section on the torus torsor has only finitely
many weights. On a finite affine cover trivializing the torsor this
is finite Laurent support, and one takes the union of the supports.
Consequently, for fixed \(n\), away from finitely many \(y\) all germs
of \(f_y\) belong to \(\pi^{m+n}\). In integral frames the preceding
equation at \(x=0\) now reads \(b f_y=a h\) up to a unit, with
\(\psi(y)=a/b\) and \(h=t_{c(y)}^*f_0\) at its witness.
If \(v_T(a/b)<n\), write \(a/b=u\pi^j\) with \(u\in T^\times\)
and \(j<n\). The local ring is \(T\)-flat, so multiplication by \(\pi\)
is injective. Clearing denominators in \(f_y=u\pi^j h\) then gives
\(h\in\pi^{m+n-j}\), a contradiction. Thus \(\{y:Q(y)<n\}\) is
finite. If \(0\ne y\) had \(Q(y)\le0\), the distinct multiples
\(ky\), \(k\ge1\), would all belong to \(\{Q<1\}\).
This proves \(Q(y)>0\) for nonzero lattice vectors. The integral
symmetric matrix of \(B\) is positive definite on the real span:
a negative direction would contain rational approximants of negative
square, and a singular rational matrix would have a rational kernel
vector. A finite extension of complete discrete valuation rings of
ramification index \(e\) sends \(Q,B\) to \(eQ,eB\) for
integer-normalized orders. This statement precedes the subsequent
use of a divisible value group on \(\overline K\).
\footnote{\href{https://www.kwlan.org/articles/cpt-PEL-type-thesis-revision-dense.pdf}{Lan,
equations (4.3.1.2), (4.3.1.7), Proposition 4.3.1.9, and
\S4.2.4}. In the release see
\srcid{ExistsDegenerationDataFunctionalEquationDegreeOn},
\srcid{PsiTendstoZero}, \srcid{CubeRelScalarsMemMaximalIdeal}
and \srcid{EvenQuadraticPosOfFiniteNonpos} in
\srcid{UniformizationProofs}.}

The polarization and the integral action of \(\mathcal O_E\) also
give the balance identity required in Section~\ref{sec:period-basis}.
Write \(\mathcal P_B\) for the Poincaré biextension. On \(B\) the
Rosati condition is \(a^\vee\lambda_B=\lambda_B\bar a\). The
defining property of the dual morphism identifies
\((a\times1)^*\mathcal P_B\) with \((1\times a^\vee)^*\mathcal P_B\).
Under this identification the induced actions on the period
trivialization satisfy
\[
        (a\times1)^*\tau=(1\times\bar a)^*\tau
                \quad\text{on }Y\times X .
\]
This is the endomorphism part of the degeneration equivalence:
the integral action on the generic polarized abelian variety
corresponds to actions on its abelian part, its two lattices and its
period map. The left action on \(X\) is obtained from its natural
right action by conjugation. The displayed equality is an equality
in integral line fibers after extending scalars to \(K\); the
canonical identification of their integral lattices comes from the
Poincaré biextension over \(T\). In independent integral frames its
coefficient is a unit. Taking valuations thus gives
\(B(ay,x)=B(y,\bar a x)\), with no scalar error. This also gives the
conjugation in the torus action used below.
\footnote{\href{https://www.kwlan.org/articles/cpt-PEL-type-thesis-revision-dense.pdf}{Lan,
Definition 5.1.1.3(4) and Theorem 5.1.1.4}. The release takes the
valuation of an origin-normalized Rosati generator in
\href{https://github.com/ChristianKudeba/alpoge-rh-verification/blob/alpoge-source-7cbf7a0c/src/ReciprocityTheory/ReductiveGroups/ArithmeticGeometry/DiscreteValuationRings/ExistsRosatiGen.lean}{\srcid{bOf_balance}}.}

\subsection{The rank bound from the PEL determinant condition}

Let \(d=[E:\mathbb Q]\), \(r=\operatorname{rank}_{\mathbb Z}X\), and
\(s=\dim_E(X\otimes_{\mathbb Z}\mathbb Q)\). For a generic point of
the moduli datum of signature \((p_0,q_0)\) we prove
\[
   r=ds,\qquad s\le\min(p_0,q_0).
\]
The first equality follows by restriction of scalars.
Put \(F=\overline K\), and fix an embedding
\(\sigma:\mathcal O_E\to F\). It suffices to exhibit \(s\) independent
\(\sigma\)-eigenvectors in \(\operatorname{Lie}(G_F)\).
We use row vectors and let \(M(a)\) be the integral matrix of \(a\)
on the Lie algebra of \(G\), in a basis over \(T\). The toric Lie
algebra at level \(e\) consists of the additive forms on \(X\) with
values in \(T/\pi^{e+1}\). This description follows by evaluating
a character on dual numbers as \(1+\varepsilon t\). In a
\(\mathbb Z\)-basis of \(X\), define \(A(\bar a)\) to be the matrix
of \(\bar a\) on \(X\) with column coordinates. The conjugation is
the source's convention for the toric part of the endomorphism.

The kernel embedding on level \(e\) induces an injection on tangent
vectors. Lift its matrix to an \(r\times g\) matrix \(C\) over
\(T\), where \(g=\dim G_K\). Equivariance and injectivity give
\[
   CM(a)\equiv A(\bar a)C\pmod{\pi^{e+1}},
   \qquad CD=I_r
\]
for some \(g\times r\) matrix \(D\) over \(T\). For the second
assertion, a residue row in the kernel of \(C\bmod\pi\) lifts to
\(w\) with \(wC\in\pi T^g\). Then \(\pi^ew\) is in the level-\(e\)
kernel, so \(\pi^ew\in\pi^{e+1}T^r\) and the original residue row
is zero. Thus \(C\) has an invertible \(r\times r\) minor over \(T\),
which gives a right inverse. Neither the lifts \(C\) nor their right
inverses need to be chosen compatibly as \(e\) varies.

Let \(V\) be the integral closure of \(T\) in \(F\).
It is a domain with fraction field \(F\). Its element \(\pi\) is a
nonunit. Every \(0\ne u\in V\) divides some power of \(\pi\) in \(V\):
a monic integral equation for \(u\) can be chosen with nonzero
constant term, giving \(u\mid b\) for some \(0\ne b\in T\);
then \(b\) is a unit times a power of \(\pi\). Also
\(\sigma(\mathcal O_E)\subset V\). In an \(E\)-basis of \(X_{\mathbb Q}\),
applying \(\sigma\circ\text{conjugation}\) to each coordinate gives \(s\)
independent additive forms with
\(\varphi(\bar a x)=\sigma(a)\varphi(x)\). Clearing a common
denominator on the integral basis of \(X\) gives an \(s\times r\)
matrix \(\Phi\) over \(V\) of full row rank satisfying
\(\Phi A(\bar a)=\sigma(a)\Phi\). Form the matrix \(H\) by adjoining
horizontally the blocks \(M(a)-\sigma(a)I_g\) for integral additive
generators \(a\) of \(\mathcal O_E\). At every level the previous
congruence gives matrices over \(V\) such that
\[
    CD=I_r,\qquad \Phi C H=\pi^{e+1}Z .
\]

Over \(F\) choose a generalized inverse of \(H\) by inverting its
restriction to a complement of its kernel on its image. Choose a
right inverse to \(\Phi\). Clearing denominators gives integral
matrices \(P,\Psi\) and \(0\ne b,b'\in V\) with
\(HPH=bH\) and \(\Phi\Psi=b'I_s\). Take \(e\) with
\(\pi^e=bb' u\), and use matrices \(C,D,Z\) for this \(e\). Define
\[
   W=b\Phi C-\Phi CH P,\qquad N=D\Psi .
\]
Then
\[
  WH=0,\qquad WN=bb'\bigl(I_s-\pi u ZP D\Psi\bigr).
\]
The determinant in parentheses is \(1\bmod \pi\). Since
\(V/(\pi)\ne0\), it is nonzero. Thus the \(s\) rows of \(W\) are
independent over \(F\). The identity \(WH=0\) puts each row in the
\(\sigma\)-eigenspace for all additive generators, hence for every
element of \(\mathcal O_E\). Finally choose \(\sigma\) and its
conjugate. The distinguished imaginary generator of the host maps
under these embeddings to the two opposite nonzero roots in the
determinant condition. Their eigenspaces have dimensions \(p_0,q_0\)
in some order. Applying the bound to both gives the assertion.
\footnote{The source proof is
\href{https://github.com/ChristianKudeba/alpoge-rh-verification/blob/alpoge-source-7cbf7a0c/src/Development/AbelianVarietyTheory/Proofs/UnitaryExistence/ExistsMulEqSmulOne.lean}{\srcid{ExistsMulEqSmulOne}},
especially \srcid{exists_matrix_level}, \srcid{exists_solutions},
\srcid{finrank_le_of_lieMult}, and \srcid{rank_le_signature}.
The argument uses the actual \(\mathcal O_E\)-equivariance on the
formal trait and the generic moduli determinant condition. The
trace-normalized pullback of the theta line and its integral torsion
identification are additional structures used in the jet calculation
of Sections~\ref{sec:theta-lift-valuation}
and~\ref{sec:local-theta-jets}.}
